\documentclass[12pt]{article}
\usepackage{amsmath, amssymb}
\usepackage[margin=1in]{geometry}
\setlength{\parskip}{6pt}
\title{QUBO Formulation: Mutual-Information Sensor Placement Problem}
\date{}
\begin{document}
\maketitle

\subsection*{Mutual-Information Sensor Placement Problem}

\paragraph{Problem Statement.}
Given a set of $N$ candidate sensor locations, the goal is to select exactly $K$ locations to maximize the total information score while minimizing pairwise measurement redundancy. Each candidate location $i$ is associated with an individual information score $I_i$, and each distinct pair of locations $(i,j)$ is associated with a redundancy penalty $R_{ij}$ that is incurred when both locations are simultaneously selected. The task is to choose $K$ sensor positions such that the sum of their individual information scores minus the sum of their pairwise redundancy penalties is maximized.

\paragraph{Decision Variables.}
For each candidate location $i \in \{0, 1, \dots, N-1\}$, define a binary decision variable $x_i \in \{0,1\}$. The variable $x_i = 1$ indicates that location $i$ is selected, and $x_i = 0$ indicates that it is not. The vector $x = (x_0, x_1, \dots, x_{N-1}) \in \{0,1\}^N$ fully specifies a candidate placement.

\paragraph{QUBO Derivation.}
\textbf{Step 1.} The original problem is a maximization. Standard QUBO solvers minimize an energy function $H(x)$. We negate the objective to convert it to a minimization problem:
\begin{align}
\min_{x \in \{0,1\}^N} \left( -\sum_{i=0}^{N-1} I_i x_i + \sum_{0 \leq i < j \leq N-1} R_{ij} x_i x_j \right).
\end{align}

\textbf{Step 2.} The selection constraint requires exactly $K$ sensors to be chosen:
\begin{align}
\sum_{i=0}^{N-1} x_i = K.
\end{align}

\textbf{Step 3.} We enforce this constraint using a quadratic penalty term with weight $\Lambda > 0$:
\begin{align}
P(x) &= \Lambda \left( \sum_{i=0}^{N-1} x_i - K \right)^2. \tag{a}
\end{align}
Expanding the square and applying the idempotent property $x_i^2 = x_i$ for binary variables yields:
\begin{align}
P(x) &= \Lambda \left( \sum_{i=0}^{N-1} x_i^2 - 2K \sum_{i=0}^{N-1} x_i + K^2 \right) \nonumber \\
&= \Lambda \left( \sum_{i=0}^{N-1} x_i - 2K \sum_{i=0}^{N-1} x_i + K^2 + \sum_{i \neq j} x_i x_j \right) \nonumber \\
&= \Lambda \left( (1 - 2K) \sum_{i=0}^{N-1} x_i + 2 \sum_{0 \leq i < j \leq N-1} x_i x_j + K^2 \right). \tag{b}
\end{align}
The resulting general penalty contribution to the Hamiltonian is:
\begin{align}
P(x) = \Lambda(1 - 2K) \sum_{i=0}^{N-1} x_i + 2\Lambda \sum_{0 \leq i < j \leq N-1} x_i x_j + \Lambda K^2. \tag{c}
\end{align}

\textbf{Step 4.} The total QUBO Hamiltonian $H(x)$ is the sum of the negated objective and the penalty term:
\begin{align}
H(x) &= \left( -\sum_{i=0}^{N-1} I_i x_i + \sum_{0 \leq i < j \leq N-1} R_{ij} x_i x_j \right) + \left( \Lambda(1 - 2K) \sum_{i=0}^{N-1} x_i + 2\Lambda \sum_{0 \leq i < j \leq N-1} x_i x_j + \Lambda K^2 \right).
\end{align}

\textbf{Step 5.} Collecting linear and quadratic terms yields the standard QUBO form $H(x) = \sum_{i} Q_{ii} x_i + \sum_{i < j} Q_{ij} x_i x_j + c$. The matrix entries and offset are given in closed form by:
\begin{align}
Q_{ii} &= -I_i + \Lambda(1 - 2K), \quad \text{for all } i \in \{0, \dots, N-1\}, \\
Q_{ij} &= R_{ij} + 2\Lambda, \quad \text{for all } 0 \leq i < j \leq N-1, \\
c &= \Lambda K^2.
\end{align}

\paragraph{Penalty Weight Justification.}
The penalty weight $\Lambda$ must be chosen sufficiently large to ensure that any violation of the cardinality constraint increases the total energy more than any possible improvement in the original objective. The maximum possible value of the original objective is bounded above by $\sum_{i=0}^{N-1} I_i + \sum_{0 \leq i < j \leq N-1} R_{ij}$. Therefore, selecting $\Lambda > \sum_{i=0}^{N-1} I_i + \sum_{0 \leq i < j \leq N-1} R_{ij}$ guarantees that the penalty term strictly dominates any feasible objective gain, forcing the global minimum to satisfy $\sum_i x_i = K$.

\paragraph{Correctness Argument.}
Minimizing $H(x)$ over $\{0,1\}^N$ recovers the optimal sensor placement because feasible configurations (those satisfying $\sum_i x_i = K$) incur zero penalty and directly minimize the negated information-redundancy objective. Any infeasible configuration violates the cardinality constraint, incurring a penalty strictly greater than the maximum possible reduction in the original objective. Consequently, the energy of every infeasible solution exceeds that of the best feasible solution, ensuring that the global minimum of $H(x)$ corresponds exactly to a feasible, optimal selection of $K$ sensor locations.

\paragraph{Illustrative Example.}
Consider a small instance with $N=3$ candidate locations and a target selection of $K=2$. The parameters are the individual scores $I_0, I_1, I_2$ and the symmetric redundancy penalties $R_{01}, R_{02}, R_{12}$. Substituting $N=3$ and $K=2$ into the general formulas yields the explicit Hamiltonian:
\begin{align}
H(x) &= \sum_{i=0}^{2} \left( -I_i + \Lambda(1 - 4) \right) x_i + \sum_{0 \leq i < j \leq 2} \left( R_{ij} + 2\Lambda \right) x_i x_j + 4\Lambda \nonumber \\
&= \sum_{i=0}^{2} \left( -I_i - 3\Lambda \right) x_i + \left( R_{01} + 2\Lambda \right) x_0 x_1 + \left( R_{02} + 2\Lambda \right) x_0 x_2 + \left( R_{12} + 2\Lambda \right) x_1 x_2 + 4\Lambda.
\end{align}
The corresponding QUBO coefficients are $Q_{00} = -I_0 - 3\Lambda$, $Q_{11} = -I_1 - 3\Lambda$, $Q_{22} = -I_2 - 3\Lambda$, $Q_{01} = R_{01} + 2\Lambda$, $Q_{02} = R_{02} + 2\Lambda$, $Q_{12} = R_{12} + 2\Lambda$, and offset $c = 4\Lambda$. For a concrete assignment $I = [10, 8, 5]$, $R_{01}=2$, $R_{02}=3$, $R_{12}=1$, and $\Lambda = 25$ (satisfying the dominance condition), the three feasible configurations and their energies are:
\begin{align}
x = (1,1,0) &\implies H(x) = -18 + 2 + 4(25) = -16 + 100 = 84, \\
x = (1,0,1) &\implies H(x) = -15 + 3 + 4(25) = -12 + 100 = 88, \\
x = (0,1,1) &\implies H(x) = -13 + 1 + 4(25) = -12 + 100 = 88.
\end{align}
The optimal configuration is $x^* = (1,1,0)$, which selects locations $0$ and $1$, yielding a maximum total information score of $10 + 8 - 2 = 16$ and a minimum energy $H(x^*) = 84$.
\end{document}